% GENERATED FILE. Edit canonical lessons, then run npm run generate:books.
% canonical-source-hash: fnv1a64:302192c25539397d
% canonical-lessons: GU-C188-maphat, GU-C188-nakli, GU-C188-asli, GU-C188-mal, GU-C188-vastu, GU-R188-review-183-188

\chapter{Genuine goods and fakes}
\label{ch:gu-genuine-goods-and-fakes}
\hlchaptermodality{188}

\begin{chapteropening}
\textbf{By the end of this chapter:} \emph{I can say whether goods are free, fake or genuine.}

Complete the last lesson of chapter 188: I can say whether goods are free, fake or genuine.
\end{chapteropening}

\section[\texorpdfstring{\gu{મફત} (maphat)}{મફત (maphat)}]{\gu{મફત} (maphat) — free of charge}

\label{lesson:GU-C188-maphat}



\begin{quote}
Before the new one: say the Gujarati for per cent, then the Gujarati for discount.
\end{quote}

\subsection*{You'll want to know: \gu{મફત}}
\textbf{\gu{મફત}} — \emph{maphat} — \textquotedblleft{}free of charge\textquotedblright{}.

\begin{quote}
\textbf{\gu{અહીં} \gu{પાણી} \gu{મફત} \gu{છે}.} — \emph{ahī̃ pāṇī maphat che} — \textquotedblleft{}Water is free here.\textquotedblright{}
\end{quote}

\begin{cousinweb}
\textbf{\gu{મફત}} (\emph{maphat}) is a loan, usually traced to Persian \emph{muft}, \textquotedblleft{}free, gratis\textquotedblright{}. CDIAL has no entry for it.
\end{cousinweb}

\begin{grammarlens}[title={What you've built}]
One more word for everyday life at A2.
\end{grammarlens}

\subsection*{Your turn}
\begin{itemize}
\raggedright
  \item \emph{Say it:} \emph{maphat}
  \item \emph{Say it:} \emph{maphat}, once more
  \item \emph{Say it:} \emph{ahī̃ pāṇī maphat che} — Water is free here.
  \item \emph{Recall:} say the Gujarati for per cent, then the Gujarati for discount, then say \emph{maphat} again
\end{itemize}

\begin{tcolorbox}[breakable,colback=teal!4,colframe=teal!35!black,title={Before you move on}]
What does \gu{મફત} mean? (\textquotedblleft{}Free of charge\textquotedblright{}.) Say it once more.
\end{tcolorbox}
\section[\texorpdfstring{\gu{નકલી} (naklī)}{નકલી (naklī)}]{\gu{નકલી} (naklī) — fake, counterfeit}

\label{lesson:GU-C188-nakli}



\begin{quote}
Before the new one: say the Gujarati for discount, then the Gujarati for free of charge.
\end{quote}

\subsection*{You'll want to know: \gu{નકલી}}
\textbf{\gu{નકલી}} — \emph{naklī} — \textquotedblleft{}fake, counterfeit\textquotedblright{}.

\begin{quote}
\textbf{\gu{આ} \gu{નોટ} \gu{નકલી} \gu{છે}.} — \emph{ā noṭ naklī che} — \textquotedblleft{}This note is fake.\textquotedblright{}
\end{quote}

\begin{cousinweb}
\textbf{\gu{નકલી}} (\emph{naklī}) is built on \emph{nakal}, \textquotedblleft{}a copy\textquotedblright{}, a loan from Arabic \emph{naql}, \textquotedblleft{}transfer, copying\textquotedblright{}, which reached Gujarati through Persian; the ending \emph{-ī} makes the adjective. Its Arabic root is \emph{n-q-l} (Lane).
\end{cousinweb}

\begin{grammarlens}[title={What you've built}]
One more word for everyday life at A2.
\end{grammarlens}

\subsection*{Your turn}
\begin{itemize}
\raggedright
  \item \emph{Say it:} \emph{naklī}
  \item \emph{Say it:} \emph{naklī}, once more
  \item \emph{Say it:} \emph{ā noṭ naklī che} — This note is fake.
  \item \emph{Recall:} say the Gujarati for discount, then the Gujarati for free of charge, then say \emph{naklī} again
\end{itemize}

\begin{tcolorbox}[breakable,colback=teal!4,colframe=teal!35!black,title={Before you move on}]
What does \gu{નકલી} mean? (\textquotedblleft{}Fake, counterfeit\textquotedblright{}.) Say it once more.
\end{tcolorbox}
\section[\texorpdfstring{\gu{અસલી} (aslī)}{અસલી (aslī)}]{\gu{અસલી} (aslī) — genuine, real}

\label{lesson:GU-C188-asli}



\begin{quote}
Before the new one: say the Gujarati for free of charge, then the Gujarati for fake.
\end{quote}

\subsection*{You'll want to know: \gu{અસલી}}
\textbf{\gu{અસલી}} — \emph{aslī} — \textquotedblleft{}genuine, real\textquotedblright{}.

\begin{quote}
\textbf{\gu{શું} \gu{આ} \gu{નોટ} \gu{અસલી} \gu{છે}?} — \emph{śũ ā noṭ aslī che?} — \textquotedblleft{}Is this note genuine?\textquotedblright{}
\end{quote}

\begin{cousinweb}
\textbf{\gu{અસલી}} (\emph{aslī}) is built on \emph{asal}, \textquotedblleft{}original\textquotedblright{}, a loan from Arabic \emph{aṣl}, \textquotedblleft{}root, origin\textquotedblright{}, which reached Gujarati through Persian; the ending \emph{-ī} makes the adjective. Its Arabic root is \emph{ʾ-ṣ-l} (Lane).
\end{cousinweb}

\begin{grammarlens}[title={What you've built}]
One more word for everyday life at A2.
\end{grammarlens}

\subsection*{Your turn}
\begin{itemize}
\raggedright
  \item \emph{Say it:} \emph{aslī}
  \item \emph{Say it:} \emph{aslī}, once more
  \item \emph{Say it:} \emph{śũ ā noṭ aslī che?} — Is this note genuine?
  \item \emph{Recall:} say the Gujarati for free of charge, then the Gujarati for fake, then say \emph{aslī} again
\end{itemize}

\begin{tcolorbox}[breakable,colback=teal!4,colframe=teal!35!black,title={Before you move on}]
What does \gu{અસલી} mean? (\textquotedblleft{}Genuine, real\textquotedblright{}.) Say it once more.
\end{tcolorbox}
\section[\texorpdfstring{\gu{માલ} (māl)}{માલ (māl)}]{\gu{માલ} (māl) — goods, stock}

\label{lesson:GU-C188-mal}



\begin{quote}
Before the new one: say the Gujarati for fake, then the Gujarati for genuine.
\end{quote}

\subsection*{You'll want to know: \gu{માલ}}
\textbf{\gu{માલ}} — \emph{māl} — \textquotedblleft{}goods, stock\textquotedblright{}.

\begin{quote}
\textbf{\gu{દુકાનમાં} \gu{નવો} \gu{માલ} \gu{આવ્યો}.} — \emph{dukānmā̃ navo māl āvyo} — \textquotedblleft{}New stock has come into the shop.\textquotedblright{}
\end{quote}

\begin{cousinweb}
\textbf{\gu{માલ}} (\emph{māl}) is a loan from Arabic \emph{māl}, \textquotedblleft{}property, wealth, whatever one owns\textquotedblright{}, which reached Gujarati through Persian. Its Arabic root is \emph{m-w-l} (Lane).
\end{cousinweb}

\begin{grammarlens}[title={What you've built}]
One more word for everyday life at A2.
\end{grammarlens}

\subsection*{Your turn}
\begin{itemize}
\raggedright
  \item \emph{Say it:} \emph{māl}
  \item \emph{Say it:} \emph{māl}, once more
  \item \emph{Say it:} \emph{dukānmā̃ navo māl āvyo} — New stock has come into the shop.
  \item \emph{Recall:} say the Gujarati for fake, then the Gujarati for genuine, then say \emph{māl} again
\end{itemize}

\begin{tcolorbox}[breakable,colback=teal!4,colframe=teal!35!black,title={Before you move on}]
What does \gu{માલ} mean? (\textquotedblleft{}Goods, stock\textquotedblright{}.) Say it once more.
\end{tcolorbox}
\section[\texorpdfstring{\gu{વસ્તુ} (vastu)}{વસ્તુ (vastu)}]{\gu{વસ્તુ} (vastu) — thing, item}

\label{lesson:GU-C188-vastu}



\begin{quote}
Before the new one: say the Gujarati for genuine, then the Gujarati for goods.
\end{quote}

\subsection*{You'll want to know: \gu{વસ્તુ}}
\textbf{\gu{વસ્તુ}} — \emph{vastu} — \textquotedblleft{}thing, item\textquotedblright{}.

\begin{quote}
\textbf{\gu{આ} \gu{દુકાનમાં} \gu{દરેક} \gu{વસ્તુ} \gu{સસ્તી} \gu{છે}.} — \emph{ā dukānmā̃ darek vastu sastī che} — \textquotedblleft{}Everything in this shop is cheap.\textquotedblright{}
\end{quote}

\begin{cousinweb}
\textbf{\gu{વસ્તુ}} (\emph{vastu}) is Sanskrit \emph{vástu}, \textquotedblleft{}existing thing; affair\textquotedblright{}, taken over whole as a learned loan, a tatsama (CDIAL 11449).
\end{cousinweb}

\begin{grammarlens}[title={What you've built}]
One more word for everyday life at A2.
\end{grammarlens}

\subsection*{Your turn}
\begin{itemize}
\raggedright
  \item \emph{Say it:} \emph{vastu}
  \item \emph{Say it:} \emph{vastu}, once more
  \item \emph{Say it:} \emph{ā dukānmā̃ darek vastu sastī che} — Everything in this shop is cheap.
  \item \emph{Recall:} say the Gujarati for genuine, then the Gujarati for goods, then say \emph{vastu} again
\end{itemize}

\begin{tcolorbox}[breakable,colback=teal!4,colframe=teal!35!black,title={Before you move on}]
What does \gu{વસ્તુ} mean? (\textquotedblleft{}Thing, item\textquotedblright{}.) Say it once more.
\end{tcolorbox}
\section[(dialogue)]{Shopping and prices: a review}

\label{lesson:GU-R188-review-183-188}



\begin{quote}
Say the Gujarati for customer and shopkeeper.
\end{quote}

\subsection*{Your turn}
Say these aloud:

\begin{itemize}
\raggedright
  \item customer, shopkeeper, groceries, to weigh, kilo
  \item rate, to cost, weight, to measure, to ask for
  \item to haggle, deal, less, pair, to test
  \item to pay, cash, small change, banknote, receipt
  \item to go up, to go down, more, per cent, discount
  \item free of charge, fake, genuine, goods, thing
\end{itemize}

\begin{tcolorbox}[breakable,colback=teal!4,colframe=teal!35!black,title={Before you move on}]
Customer? (\textbf{\gu{ગ્રાહક}}.) Shopkeeper? (\textbf{\gu{દુકાનદાર}}.) Groceries? (\textbf{\gu{કરિયાણું}}.) To weigh? (\textbf{\gu{તોલવું}}.) Kilo? (\textbf{\gu{કિલો}}.) Rate? (\textbf{\gu{ભાવ}}.) To cost? (\textbf{\gu{લાગવું}}.) Weight? (\textbf{\gu{વજન}}.) To measure? (\textbf{\gu{માપવું}}.) To ask for? (\textbf{\gu{માગવું}}.) To haggle? (\textbf{\gu{ભાવતાલ} \gu{કરવો}}.) Deal? (\textbf{\gu{સોદો}}.) Less? (\textbf{\gu{ઓછું}}.) Pair? (\textbf{\gu{જોડી}}.) To test? (\textbf{\gu{પારખવું}}.) To pay? (\textbf{\gu{ચૂકવવું}}.) Cash? (\textbf{\gu{રોકડ}}.) Small change? (\textbf{\gu{પરચૂરણ}}.) Banknote? (\textbf{\gu{નોટ}}.) Receipt? (\textbf{\gu{રસીદ}}.) To go up? (\textbf{\gu{વધવું}}.) To go down? (\textbf{\gu{ઘટવું}}.) More? (\textbf{\gu{વધારે}}.) Per cent? (\textbf{\gu{ટકો}}.) Discount? (\textbf{\gu{છૂટ}}.) Free of charge? (\textbf{\gu{મફત}}.) Fake? (\textbf{\gu{નકલી}}.) Genuine? (\textbf{\gu{અસલી}}.) Goods? (\textbf{\gu{માલ}}.) Thing? (\textbf{\gu{વસ્તુ}}.)
\end{tcolorbox}
